\section{STRIDE Threat Modelling Framework}\label{sec:stride}

The STRIDE threat modeling framework helps to group real-world attacks into six types. Each type targets a different part of the authentication process. This section maps the vulnerabilities discussed earlier to STRIDE categories and shows how each of the three authentication methods handles them.

\vspace{0.3em}
\noindent
\textbf{STRIDE categories:}
\begin{itemize}
  \item \textbf{Spoofing}: when an attacker pretends to be a valid user to gain access
  \item \textbf{Tampering}: when an attacker changes messages or data during login
  \item \textbf{Repudiation}: when a user denies doing an action, like logging in
  \item \textbf{Information Disclosure}: when sensitive data like passwords or OTPs are leaked
  \item \textbf{Denial of Service}: when the attacker blocks the user from logging in
  \item \textbf{Elevation of Privilege}: when someone gets more access than they should
\end{itemize}

Each vulnerability from Section~\ref{sec:common-vulnerabilities} fits one or more of these categories. Table~\ref{tab:stride-vuln-matrix} shows which threat is mapped where, and what to expect from password-only, 2FA, and FIDO2-based systems.

\begin{table}[H]
\centering
\caption{STRIDE Categories, Real-World Vulnerabilities, and Expected Behavior}
\label{tab:stride-vuln-matrix}
\begin{tabular}{|p{2.3cm}|p{4.3cm}|p{3.3cm}|p{3.3cm}|p{3.3cm}|}
\hline
\textbf{STRIDE Category} & \textbf{Mapped Vulnerability} & \textbf{Password-only} & \textbf{2FA} & \textbf{FIDO2/WebAuthn} \\
\hline
Spoofing & Phishing and OTP relay attacks & Easy to spoof if password is stolen or phished & OTP relay attacks can still trick users & Origin check and public-key signatures prevent spoofing \\
\hline
Tampering & Adversary-in-the-middle modifying login flow & No strong protection if TLS fails & Relayed OTPs or injected prompts possible & Signed challenge-response blocks message tampering \\
\hline
Repudiation & No way to prove who logged in & Anyone with the password can log in & 2FA logs may help but are weak proof & Authenticator signs each login, tied to device and origin \\
\hline
Information Disclosure & Password reuse and data breaches & Password leaks are common and reused across sites & OTP seeds and SMS codes can be stolen & No password or shared secret stored or sent \\
\hline
Denial of Service & Push fatigue and SIM swap & Lockout by brute force or repeated failures & Flood of OTP or push requests can lock out users & Brute-force blocked by design, backup keys help recovery \\
\hline
Elevation of Privilege & Use of weak or reused admin credentials & Admin access can be guessed or reused & OTP can be phished from high-priv users & High-priv accounts tied to secure device credentials \\
\hline
\end{tabular}
\end{table}

The STRIDE mapping helps compare which attack types are stopped or remain possible in each authentication setup. This forms the basis of the analysis in Section~\ref{sec:stride-analysis}, where each system’s real threat resistance is examined in more detail.
